% !TeX root = ../tesis.tex

% -----------------------------------------------------------------------
\section{Teoría Latinoamericana}
\label{sec:teoria-latinoamericana}
% -----------------------------------------------------------------------

Para enmarcar el caso de la \cdmx{} en un contexto de crítica urbanística
regional,
es fundamental recurrir a la teoría sobre la ciudad en América Latina. La obra
compilatoria de \citet{padilla2007} proporciona el marco conceptual para
entender
que la crisis de movilidad en la periferia es síntoma de un agotamiento del
modelo
de planeación estatal y la adopción de enfoques que favorecieron el capital
inmobiliario y la privatización del espacio. La incapacidad de la planeación
urbana
para gestionar la expansión territorial generada por los grandes proyectos
viales,
como el Periférico, es el caldo de cultivo para la emergencia de la ciudad
neoliberal.

El escenario de un corredor anillar constituye, en este contexto teórico, un
caso de demostración metodológica: permite evaluar si la transición de un
enfoque radial ---propio de la era Uruchurtu--- hacia uno que incorpora
geometrías transversales reduce cuantitativamente los indicadores de
ineficiencia topológica que aquejaban a la ciudad fragmentada
\citep{franco2022}.

% -----------------------------------------------------------------------
\section{Casos Similares en Otros Países}
\label{sec:casos-similares}
% -----------------------------------------------------------------------

La implementación de corredores circulares de transporte masivo no es una idea
nueva a nivel global. Su efectividad ha sido documentada en diversas
grandes metrópolis
que enfrentaban problemas estructurales similares a los de la \cdmx{}: redes
radiales saturadas, periferias desconectadas y tiempos de traslado crecientes.

El caso más representativo es la \textbf{Gran Línea Circular del Metro de Moscú}
,
inaugurada en 2023 con 70 kilómetros de longitud. Esta línea fue concebida
explícitamente para descomprimir las líneas radiales del centro y articular la
movilidad entre zonas periféricas sin necesidad de pasar por las estaciones
centrales más congestionadas \citep{efe2023}. Su puesta en operación
demostró cómo la adopción de geometrías anillares puede cambiar cualitativamente
la eficiencia de todo el sistema de transporte de una capital.

Un segundo referente es el \textbf{\textit{Suburban Rail Loop} de Melbourne},
Australia, actualmente en planeación y construcción. Este proyecto parte de un
diagnóstico análogo al de la \cdmx{}: las alcaldías del llamado
\textit{Middle-Ring}
albergan hospitales, universidades y centros de empleo de alta demanda, pero
carecen de una conectividad lateral robusta que los articule entre sí
\citep{suburbanrailloop2021}. El proyecto busca transformar una serie de
suburbios desconectados en una red policéntrica de precintos de innovación y
empleo, mediante la inserción de un corredor ferroviario de arco periférico.

Estos referentes internacionales justifican la pertinencia del escenario
hipotético empleado en esta investigación: la incorporación de corredores
anillares en los cuatro cuadrantes del Anillo Periférico no es un experimento
sin precedente, sino una geometría probada en otros contextos metropolitanos,
lo que hace de su simulación un caso de demostración metodológica válido y
significativo.

